\chapter{课程使用方式与平台能力}

\section{课程参与者}

VeriSpecOSLab 服务于学生和课程教学团队。学生负责理解实验目标、记录设计、完成实现并解释结果。教师和助教组成课程教学团队，共同准备课程资料、规格模板、公开测试与硬件验收方法，同时维护平台配置、运行工具和课程案例。

平台以本地课程项目为中心。教学团队可以提供起始仓库，也可以让学生从空项目开始。学生在自己的仓库中完成设计、调用 AI 助手、构建内核、运行 QEMU 或开发板实验，并生成报告和提交包。教师和助教使用同一套规格、测试与运行记录复查学生结果，课程无需先部署在线服务。

\begin{table}[htbp]
  \centering
  \caption{课程参与者与主要工作}
  \begin{tabularx}{\textwidth}{>{\bfseries}p{2.8cm} X X}
    \toprule
    参与者 & 主要工作 & 平台提供的支持 \\
    \midrule
    学生 & 理解任务、亲手编写规格、审查 AI 实现、调试并解释结果 & 概念问答、只读评审、受控代码修改、构建运行、测试记录、报告与提交 \\
    教学团队 & 准备资料与测试、维护课程配置和案例、复查设计与运行结果 & 规格模板、课程知识库、统一运行步骤、完整日志和可复查提交包 \\
    \bottomrule
  \end{tabularx}
\end{table}

\section{平台如何支持一次实验}

平台覆盖学生完成一次操作系统实验所需的主要环节。每个环节都产生下一步可以直接使用的结果，课程要求、AI 修改、实际运行和最终提交由此保持一致。

\begin{longtable}{@{}p{2.6cm}p{6.1cm}p{5.8cm}@{}}
  \caption{平台能力与课程用途}\\
  \toprule
  课程环节 & 平台完成的工作 & 直接结果 \\
  \midrule
  \endfirsthead
  \toprule
  课程环节 & 平台完成的工作 & 直接结果 \\
  \midrule
  \endhead
  项目准备 & 创建课程项目并检查编译器、QEMU、模型配置和必要文件 & 可开始实验的项目目录和清晰的环境检查结果 \\
  设计与规格 & 学生先进行概念问答，再手写系统方案、模块职责、接口、目标和测试方法 & 经过 lint 与只读评审的当前任务说明 \\
  AI 协作实现 & 实现助手在限定文件范围内生成代码与多类测试，其他角色负责问答、调试、验证和评审 & 与当前任务一致的代码提交、测试或诊断结果 \\
  构建与运行 & 按课程项目声明的步骤执行构建、测试、QEMU 和开发板任务 & 实际日志、输出文件和明确的运行结论 \\
  资料与问答 & 导入教材、课程文档和参考代码，并在回答中标明来源 & 可以回到原文核对的课程问答 \\
  报告与提交 & 汇总代码版本、规格、测试和运行记录，生成报告与课程提交包 & 与实际实验过程对应的最终材料 \\
  \bottomrule
\end{longtable}

\section{教学团队如何验收}

传统实验提交通常只包含最终代码和一份学生撰写的报告。VeriSpecOSLab 让教学团队同时看到学生亲手提交的规格、AI 实际修改的文件、生成的测试、构建与验证命令、完整运行日志以及最终提交对应的代码版本。教师可以据此检查实现是否符合当前 Lab，追踪测试结论的来源，并在答辩中要求学生解释关键设计和代码。

QEMU 与开发板结果分别记录。QEMU 用于快速重复执行课程测试，开发板记录用于确认真实硬件上的启动和运行情况。最终课程评定由教学团队完成，平台负责提供统一、完整、便于复查的材料。

\section{实现原则}

平台遵循四项原则。规格由学生编写和提交。AI 生成的代码与测试需要经过实际构建和验证。每次运行结果都对应明确的代码版本和任务说明。学生、教师和助教使用同一套项目配置与记录。后续章节将分别说明规格、AI 助手、运行工具和课程案例如何实现这些原则。
